\chapter{Dériver, composer et contrôler les variations}
\section*{Exercice 14.1 -- Chaîne}
Les dérivées sont $6xe^{3x^2-1}$ et $2x/(1+x^2)$. Le second domaine est tout $\R$ car $1+x^2>0$. Dans chaque cas, la dérivée intérieure multiplie celle de la fonction extérieure.

\section*{Exercice 14.2 -- Quotient}
La dérivée vaut $(1-x^2)/(1+x^2)^2$. Elle est positive sur $]-1,1[$ et négative à l'extérieur. La fonction décroît, puis croît, puis décroît ; elle a un minimum local $-1/2$ en $-1$ et un maximum local $1/2$ en $1$.

\section*{Exercice 14.3 -- Tangente}
$\ln1=0$ et $\ln'(1)=1$, donc la tangente est $y=x-1$. On a $\ln(1+h)=h+h\varepsilon(h)$, avec $\varepsilon(h)\to0$ : l'erreur est petite devant $|h|$, sans être identiquement nulle.

\section*{Exercice 14.4 -- Dérivée nulle}
$f(x)=x^3$ convient : $f'(0)=0$ mais $f$ est strictement croissante et change de signe autour de zéro. Une annulation de dérivée n'est qu'un candidat pour un extremum intérieur.

\section*{Exercice 14.5 -- Borne d’erreur}
Si $x\ne y$, les accroissements finis donnent $\sin x-\sin y=\cos(c)(x-y)$ pour un $c$ entre eux. Comme $|\cos c|\le1$, la borne suit. Si $x=y$, les deux membres sont nuls.

\section*{Exercice 14.6 -- Optimisation}
Avec un côté $x>0$, l'autre vaut $36/x$. Le périmètre est $P(x)=2x+72/x$, de dérivée $2-72/x^2$. Il décroît jusqu'à $6$ puis croît. Le minimum est $P(6)=24$, atteint pour le carré.
